\begin{textsample}{NEG Dim 2 – pb – Score -2.00 – pb/pb\_inf\_21.txt}  \label{ex:f2_neg_008}
Campo de experiências : “ ESCUTA, FALA, PENSAMENTO E IMAGINAÇÃO ” OBJETIVOS DE APRENDIZAGEM E DESENVOLVIMENTO Dialogar com crianças e adultos, expressando seus desejos, necessidades, sentimentos e opiniões.
Identificar e criar diferentes sons e reconhecer rimas e aliterações em cantigas de roda e textos poéticos.
Demonstrar interesse e atenção ao ouvir a leitura de histórias e outros textos, diferenciando escrita de ilustrações, e acompanhando, com orientação do adulto-leitor, a direção da leitura ( de cima para baixo, da esquerda para a direita ).
Formular e responder perguntas sobre fatos da história narrada, identificando cenários, personagens e principais acontecimentos.
Relatar experiências e fatos acontecidos, histórias ouvidas, filmes ou peças teatrais assistidos etc.
Criar e contar histórias oralmente, com base em imagens ou temas sugeridos.
Manusear diferentes portadores textuais, demonstrando reconhecer seus usos sociais.
Manipular textos e participar de situações de escuta para ampliar seu contato com diferentes gêneros textuais ( parlendas, histórias de aventura, tirinhas, cartazes de sala, cardápios, notícias etc. ).
Manusear diferentes instrumentos e suportes de escrita para desenhar, traçar \textbf{letras} e outros sinais gráficos. ( BNCC, 2017 ) Vivenciar diferentes sons, rimas e aliterações com cantigas de roda e textos poéticos.
SUGESTÕES DE VIVÊNCIAS QUE FAVORECEM A CONSTRUÇÃO DE EXPERIÊNCIAS PELAS CRIANÇAS BEM PEQUENAS NESTE CAMPO As práticas realizadas na Instituição de Educação Infantil devem possibilitar : - Condições para que as crianças possam expressar sentimentos, desejos e necessidades usando a linguagem oral ; - Contato com diferentes gêneros textuais ; - Momentos de encantamento pela escuta da leitura de histórias e outros gêneros textuais ; - Participação em situações que envolvam a convivência e a manipulação de diferentes suportes textuais ; - Momentos de criação e recontos de histórias da literatura universal, local e da vida cotidiana, com ênfase nos múltiplos elementos que compõem uma narrativa ; - Participação em variadas possibilidades de criação e expressão oral, mobilizadas por materiais diversos, imagens, objetos, fotografias, filmes, \textbf{cenas} cotidianas, dentre outros recursos ; - Intervenções nas contações de histórias, de forma lúdica, utilizando múltiplas linguagens ; - Liberdade de se expressar, transmitindo suas emoções, sentimentos, pensamentos, desejos, necessidades, por meio de diferentes linguagens ; - Desenvolvimento da imaginação, atenção e percepção em suas variadas dimensões ;

% matched lemmas: cena, letra
\end{textsample}
